\documentclass[10pt,a4paper]{amsart}
\usepackage[margin=1in]{geometry}
\usepackage{barnette-figs}
\usepackage{graphicx,amssymb}
\setlength{\parindent}{0pt}
\setlength{\parskip}{7pt}
\pagestyle{plain}
\begin{document}
\section*{22. Counts separate existence from abundance}
\begin{center}\begin{minipage}{.86\textwidth}\centering
\begin{tikzpicture}[bf,x=1cm,y=1cm]
\draw[draw=bfS,line width=1pt,] (0.000,0.600) -- (1.300,0.951) -- (2.600,1.200) -- (3.900,1.393) -- (5.200,1.551);
\node[draw=bfS,fill=white,rectangle,minimum size=3pt,inner sep=0pt] at (0.0,0.6) {};
\node[draw=bfS,fill=white,rectangle,minimum size=3pt,inner sep=0pt] at (1.3,0.9509775004326936) {};
\node[draw=bfS,fill=white,rectangle,minimum size=3pt,inner sep=0pt] at (2.6,1.2) {};
\node[draw=bfS,fill=white,rectangle,minimum size=3pt,inner sep=0pt] at (3.9000000000000004,1.3931568569324173) {};
\node[draw=bfS,fill=white,rectangle,minimum size=3pt,inner sep=0pt] at (5.2,1.5509775004326936) {};
\node[anchor=west,text=bfS] at (5.5,3.5) {states};
\draw[draw=bfR,line width=1pt,dashed] (0.000,0.600) -- (1.300,0.600) -- (2.600,0.600) -- (3.900,0.600) -- (5.200,0.600);
\node[draw=bfR,fill=white,rectangle,minimum size=3pt,inner sep=0pt] at (0.0,0.6) {};
\node[draw=bfR,fill=white,rectangle,minimum size=3pt,inner sep=0pt] at (1.3,0.6) {};
\node[draw=bfR,fill=white,rectangle,minimum size=3pt,inner sep=0pt] at (2.6,0.6) {};
\node[draw=bfR,fill=white,rectangle,minimum size=3pt,inner sep=0pt] at (3.9000000000000004,0.6) {};
\node[draw=bfR,fill=white,rectangle,minimum size=3pt,inner sep=0pt] at (5.2,0.6) {};
\node[anchor=west,text=bfR] at (5.5,3.0) {ordered pairs};
\draw[draw=bfH,line width=1pt,dash dot] (0.000,1.551) -- (1.300,1.800) -- (2.600,1.993) -- (3.900,2.151) -- (5.200,2.284);
\node[draw=bfH,fill=white,rectangle,minimum size=3pt,inner sep=0pt] at (0.0,1.5509775004326936) {};
\node[draw=bfH,fill=white,rectangle,minimum size=3pt,inner sep=0pt] at (1.3,1.7999999999999998) {};
\node[draw=bfH,fill=white,rectangle,minimum size=3pt,inner sep=0pt] at (2.6,1.9931568569324172) {};
\node[draw=bfH,fill=white,rectangle,minimum size=3pt,inner sep=0pt] at (3.9000000000000004,2.1509775004326936) {};
\node[draw=bfH,fill=white,rectangle,minimum size=3pt,inner sep=0pt] at (5.2,2.2844129532345625) {};
\node[anchor=west,text=bfH] at (5.5,2.5) {Hamiltonian cycles};
\draw[->] (-.3,0)--(5.4,0) node[right] {$p$};\draw[->] (-.3,0)--(-.3,3) node[above] {$\log_2$ count};
\node[anchor=north] at (0.0,0) {4};
\node[anchor=north] at (1.3,0) {6};
\node[anchor=north] at (2.6,0) {8};
\node[anchor=north] at (3.9000000000000004,0) {10};
\node[anchor=north] at (5.2,0) {12};
\node[anchor=east] at (-.3,0.6) {1};
\node[anchor=east] at (-.3,1.2) {2};
\node[anchor=east] at (-.3,1.7999999999999998) {3};
\node[anchor=east] at (-.3,2.4) {4};
\end{tikzpicture}\par\bigskip Width-four nanotubes, direct enumeration:\par \begin{tabular}{ccc}\hline
layers $L$ & states & tree matchings\\\hline
0 & 2 & 2\\
1 & 5 & 2\\
2 & 13 & 2\\
3 & 34 & 2\\
\hline\end{tabular}
\end{minipage}\end{center}
The chart recomputes even prisms with rim length four through twelve. The table directly enumerates width-four nanotubes through three layers. The vertical axis is logarithmic to base two. State counts grow with the rim, the ordered-pair count remains two, and Hamiltonian-cycle counts grow faster than the recovery image. The prism proofs in the family report explain these observations by explicit states. Hexagonal nanotubes behave differently: their transfer matrices remember live matching bits and connectivity partitions along a boundary. Width four keeps only two tree matchings while its state count grows; width six has growing tree-matching counts. Those nanotube tables and their certified recurrences are still unrefereed, so larger-layer claims are excluded from the comparison. A small rooted image can prove existence while becoming a poor model of a typical Hamiltonian cycle. Counting routes must specify which object they count.

\textbf{Status.} COMPUTED prisms p=4,6,8,10,12 and width-four nanotubes L=0,1,2,3 here. General prism formulas and nanotube transfers are PROVED-HERE in an unrefereed family report; no independent review claimed.

\textbf{Connects to.} 20 produced versus missed; 26 counting conjectures.

{\footnotesize\raggedright Sources: research/\allowbreak{}2026-11-14\_sol\_counting-families/\allowbreak{}FAMILIES.md;\allowbreak{} build.py;\allowbreak{} prism-counts.json.\par}

\end{document}